% ------------------------------------------------------------------------
% RAIO-X DO ECOSSISTEMA OPENCODE v5.0.0
% Relatório Técnico Abrangente para Avaliação em Banca CNPq
% Formato: ABNT NBR 6022:2018 (Artigo em Publicação Técnico-Científica)
% Compilação: pdflatex -> bibtex -> pdflatex -> pdflatex
% ------------------------------------------------------------------------
\documentclass[
  brazil,
  sumario=tradicional,
  10pt
]{abntex2}

% --- Pacotes fundamentais ---
\usepackage[utf8]{inputenc}
\usepackage{lmodern}
\usepackage[T1]{fontenc}
\usepackage{microtype}
\usepackage{amsmath,amssymb}
\usepackage{graphicx}
\usepackage{xcolor}
\usepackage{hyperref}
\usepackage{booktabs}
\usepackage{longtable}
\usepackage{listings}
\usepackage{tikz}
\usepackage[alf,abnt-etal-text=it,abnt-etal-list=0]{abntex2cite}

\usetikzlibrary{shapes.geometric, arrows.meta, positioning, shadows, fit}

% --- Configurações de cores ---
\definecolor{blueopencode}{RGB}{37,99,235}
\definecolor{greenopencode}{RGB}{22,163,74}
\definecolor{orangeopencode}{RGB}{234,88,12}
\definecolor{redopencode}{RGB}{220,38,38}
\definecolor{purpleopencode}{RGB}{124,58,237}
\definecolor{graycode}{RGB}{30,41,59}
\definecolor{lightgray}{RGB}{241,245,249}

% --- Configuração de listings ---
\lstdefinestyle{javastyle}{
  language=Java,
  basicstyle=\small\ttfamily,
  keywordstyle=\color{blueopencode}\bfseries,
  commentstyle=\color{green!50!black},
  stringstyle=\color{redopencode},
  numbers=left,
  numberstyle=\tiny,
  frame=single,
  breaklines=true,
  backgroundcolor=\color{lightgray},
  tabsize=2
}

\lstdefinestyle{pythonstyle}{
  language=Python,
  basicstyle=\small\ttfamily,
  keywordstyle=\color{purpleopencode}\bfseries,
  commentstyle=\color{green!50!black},
  stringstyle=\color{orangeopencode},
  numbers=left,
  numberstyle=\tiny,
  frame=single,
  breaklines=true,
  backgroundcolor=\color{lightgray},
  tabsize=2
}

% --- Metadados ABNT ---
\titulo{Ecossistema OpenCode v5.0.0: Uma Arquitetura Integrada de Engenharia de Software com Agentes Inteligentes, Computação Quântica e Validação Acadêmica Multiagente}
\autor{Equipe OpenCode}
\local{Brasil}
\data{Junho 2026}

\instituicao{%
  OpenCode Ecosystem
  \par
  Engenharia de Software com Agentes Inteligentes
}

\tipotrabalho{Relatório Técnico}

% ======================================================================
\begin{document}
% ======================================================================

\resumo
  Este relatório técnico apresenta o raio-x completo do ecossistema
  OpenCode v5.0.0, uma plataforma integrada de engenharia de software
  com agentes inteligentes que compreende mais de 600 componentes
  distribuídos em 125 agentes especializados, 150 habilidades (\textit{skills}),
  46 servidores MCP (\textit{Model Context Protocol}), 4 motores de
  raciocínio formal (Z3, SymPy, miniKanren, Critical), 38 \textit{skills}
  científicas integradas a bases como AlphaFold, PubMed e UniProt, e um
  subsistema de computação quântica com 107 arquivos incluindo QML
  (Quantum Machine Learning) com acurácia de 89,52\% no conjunto HAM10000.
  São documentados: (i) a arquitetura em 3 camadas (MCP $\rightarrow$
  Skill $\rightarrow$ Agente); (ii) o \textit{pipeline} acadêmico com
  49 agentes MASWOS e 91 artefatos de criação de artigos; (iii) o
  sistema de orquestração Nexus com 488 arquivos e 6 camadas de
  meta-granularidade (L0--L6); (iv) o sistema evolutivo com 14 gerações
  e evolução de score 85 para 98; (v) os 4 motores de raciocínio com
  212+ tipos em 27 categorias; (vi) o \textit{benchmark} CORA-Eval com
  150 tarefas em 10 dimensões; e (vii) o \textit{pipeline} SDD--TDD--
  Harness--Hooks com 25 assertivas críticas. O ecossistema é analisado
  sob as lentes da Engenharia de Software, Inteligência Artificial,
  Computação Quântica e Metodologia Científica, com ênfase em
  rastreabilidade, validação formal e produção acadêmica Qualis A1.
\endresumo

% ======================================================================
\section{INTRODUÇÃO}
% ======================================================================

A convergência entre engenharia de software, inteligência artificial e
computação quântica tem impulsionado o desenvolvimento de plataformas
cada vez mais integradas e autônomas \cite{sommerville2011engenharia,
pressman2016engenharia}. O ecossistema OpenCode v5.0.0 surge como uma
resposta a este cenário, propondo uma arquitetura unificada que combina
agentes inteligentes, validação formal, computação quântica e produção
acadêmica em um único ambiente de desenvolvimento.

Este relatório apresenta um raio-x completo do ecossistema, cobrindo:

\begin{enumerate}
  \item \textbf{Arquitetura em 3 camadas:} MCPs (46 servidores),
    Skills (150 habilidades), Agentes (125 especialistas).
  \item \textbf{Computação Quântica:} 107 arquivos, QML com acurácia
    de 89,52\% no HAM10000, mitigação de erros ZNE/PEC/DD.
  \item \textbf{Orquestração Multiagente Nexus:} 488 arquivos,
    6 camadas de meta-granularidade, 120+ barreiras de sincronização.
  \item \textbf{Pipeline de Criação de Artigos:} 91 arquivos,
    49 agentes MASWOS, 14 referências Qualis A1.
  \item \textbf{Sistema Evolutivo:} 14 gerações, score 85 $\rightarrow$ 98.
  \item \textbf{Motores de Raciocínio:} Z3, SymPy, miniKanren, Critical.
  \item \textbf{Skills Científicas:} 38 integrações com bases como
    AlphaFold, PubMed, ChEMBL, UniProt.
  \item \textbf{Benchmark CORA-Eval:} 150 tarefas em 10 dimensões.
  \item \textbf{Pipeline SDD--TDD--Harness--Hooks:} 25 assertivas
    críticas, validação hierárquica.
\end{enumerate}

A Seção~\ref{sec:arquitetura} descreve a arquitetura geral do ecossistema.
A Seção~\ref{sec:quantum} detalha o subsistema de computação quântica.
A Seção~\ref{sec:nexus} apresenta a orquestração multiagente Nexus.
A Seção~\ref{sec:artigo} documenta o pipeline de criação de artigos.
A Seção~\ref{sec:validadores} descreve o sistema de validadores e o
pipeline SDD--TDD--Harness--Hooks. A Seção~\ref{sec:evolucao} apresenta
o sistema evolutivo. A Seção~\ref{sec:raciocinio} documenta os motores
de raciocínio. A Seção~\ref{sec:ciencia} detalha as \textit{skills}
científicas. A Seção~\ref{sec:coraeval} apresenta o \textit{benchmark}
CORA-Eval. A Seção~\ref{sec:conclusao} tece as considerações finais.

% ======================================================================
\section{ARQUITETURA DO ECOSSISTEMA}
\label{sec:arquitetura}
% ======================================================================

A arquitetura do OpenCode v5.0.0 é organizada em 3 camadas principais,
seguindo o padrão de \textit{layered architecture}
\cite{martin2017clean}:

% ----------------------------------------------------------------------
\subsection{Camada 1: MCPs (Model Context Protocol)}
% ----------------------------------------------------------------------

A camada de MCPs constitui a interface entre o ecossistema e fontes
externas de dados e serviços. São 46 servidores MCP, dos quais 23
encontram-se ativos, distribuídos nas seguintes categorias funcionais:

\begin{longtable}{p{3cm}p{3.5cm}p{5.5cm}}
  \caption{Categorias de MCPs no ecossistema OpenCode}
  \label{tab:mcps}\\
  \toprule
  \textbf{Categoria} & \textbf{Quantidade} & \textbf{Exemplos} \\
  \midrule
  \endfirsthead
  \toprule
  \textbf{Categoria} & \textbf{Quantidade} & \textbf{Exemplos} \\
  \midrule
  \endhead
  \bottomrule
  \endfoot
  Busca             & 4  & websearch, gh\_grep, context7, scihub \\
  Navegador         & 2  & playwright, chrome-devtools \\
  Código            & 3  & eslint, diff, code-runner \\
  Dados             & 3  & sqlite, fetch, pdf, time \\
  Raciocínio        & 2  & sequential-thinking, memory \\
  Infraestrutura    & 2  & filesystem, github \\
  Ciência           & 4  & arxiv-mcp, latest-science, research-mcp, sura-papers \\
  Outros            & 3  & node-sandbox, websearch, web-fetch \\
  \hline
  \textbf{Total}    & \textbf{23 ativos} & (de 46 instalados) \\
\end{longtable}

% ----------------------------------------------------------------------
\subsection{Camada 2: Skills (Habilidades)}
% ----------------------------------------------------------------------

A camada de Skills é composta por 150 habilidades distribuídas em
13 categorias:

\begin{longtable}{p{3cm}p{2cm}p{7cm}}
  \caption{Categorias de Skills no ecossistema OpenCode}
  \label{tab:skills}\\
  \toprule
  \textbf{Categoria} & \textbf{Qtde.} & \textbf{Descrição} \\
  \midrule
  \endfirsthead
  \toprule
  \textbf{Categoria} & \textbf{Qtde.} & \textbf{Descrição} \\
  \midrule
  \endhead
  \bottomrule
  \endfoot
  System              & 12   & Auditoria acadêmica, revisão de código, \
                                      planos, MCP discovery \\
  Jurídico            & 7    & Contratos, triagem, peças jurídicas, \
                                      edição cirúrgica \\
  Research            & 18   & Editais BR, SEER, arteterapia clínica, \
                                      ML acadêmico \\
  Science             & 38   & AlphaFold, PubMed, ChEMBL, UniProt, \
                                      ClinVar, FoldSeek, PyMOL \\
  Reasoning           & 4    & Z3, SymPy, miniKanren, Critical \\
  Agency/Academic     & 5    & Anthropologist, Geographer, Historian, \
                                      Narratologist, Psychologist \\
  Open Design         & 75+  & Templates HTML, decks, dashboards, \
                                      landing pages, vídeos \\
  Broomva             & 30+  & Content engine, saúde, finanças, \
                                      oncologia, genômica \\
  Baoyu               & 20+  & Imagens, tradução, compressão, \
                                      diagramas, WeChat, X \\
  CC-Skills           & 10+  & SEO, ghostwriting, press release, \
                                      chrome-extension \\
  Agent Skills        & 20+  & Code review, TDD, debugging, CI/CD, \
                                      documentation \\
  Agente-Forum        & 10+  & MiroFish, BettaFish, ANP, OASIS, \
                                      PhD Auditor \\
  Especificações      & 10+  & CORA-Eval D1--D11, TDDAcademic \\
\end{longtable}

% ----------------------------------------------------------------------
\subsection{Camada 3: Agentes}
% ----------------------------------------------------------------------

O OpenCode dispõe de 125 agentes especializados, distribuídos em:

\begin{itemize}
  \item \textbf{56 agentes core:} Autoevolve (evolução autônoma),
    Reversa (25 subagentes de engenharia reversa), Coder/Frontend
    (8 especialistas), DevOps, Test, Security, Audit, Build.
  \item \textbf{49 agentes de criação:} Article Creator (MASWOS),
    44 agentes especializados por domínio científico mais 5 de
    suporte (scheduler, revisor, consultor, corretor, avaliador).
  \item \textbf{12 agentes SEER:} Pesquisa acadêmica com 10+
    fontes (arXiv, OpenAlex, Semantic Scholar, PubMed, CORE).
  \item \textbf{8 agentes Reversa:} Scout, Archaeologist, Architect,
    Writer, Reviewer, Synthesis, Swarm, Agent-Forum.
\end{itemize}

% ----------------------------------------------------------------------
\subsection{Integração entre Camadas}
% ----------------------------------------------------------------------

A integração entre as camadas segue o princípio de injeção de
dependência e desacoplamento via MCP. Um agente não acessa diretamente
uma fonte de dados; ele o faz através de uma Skill que, por sua vez,
utiliza um MCP específico. Este padrão arquitetural permite:

\begin{itemize}
  \item Substituição de implementações sem impacto nos agentes.
  \item Versionamento independente de MCPs, Skills e Agentes.
  \item Auditoria e rastreabilidade de todas as chamadas externas.
  \item Paralelismo na execução de múltiplos MCPs simultâneos.
\end{itemize}

A Figura~\ref{fig:arquitetura} (disponível no arquivo
\texttt{diagrams/arquitetura\_ecossistema\_opencode.svg}) ilustra a
disposição das camadas e seus componentes.

% ======================================================================
\section{SUBSISTEMA DE COMPUTAÇÃO QUÂNTICA}
\label{sec:quantum}
% ======================================================================

O subsistema quântico do OpenCode, localizado no diretório
\texttt{quantum/}, compreende 107 arquivos organizados em 5 frameworks
principais:

\begin{itemize}
  \item \textbf{Quantum Machine Learning (QML):} Implementação de
    redes neurais quânticas aplicadas à classificação de imagens
    médicas, utilizando o conjunto de dados HAM10000 (dermatoscopia).
    Acurácia alcançada: 89,52\%.
  \item \textbf{MPS (Matrix Product States):} Simulação de sistemas
    quânticos com até 50 qubits utilizando representação tensorial
    compacta (estados produto de matrizes).
  \item \textbf{Mitigação de Erros:} Pipeline ZNE (Zero-Noise
    Extrapolation), PEC (Probabilistic Error Cancellation) e DD
    (Dynamic Decoupling) para correção de erros em processadores
    quânticos NISQ (Noisy Intermediate-Scale Quantum).
  \item \textbf{Grad-CAM Quântico:} Adaptação do Grad-CAM
    (Gradient-weighted Class Activation Mapping) para interpretabilidade
    de decisões em modelos quânticos.
  \item \textbf{Referências Qualis A1:} 21 citações acadêmicas
    indexadas no sistema de classificação Qualis A1 da CAPES.
\end{itemize}

O pipeline de mitigação de erros implementa três técnicas
complementares:

\begin{enumerate}
  \item \textbf{ZNE:} Executa o mesmo circuito em diferentes níveis
    de ruído (via pulos de gate ou inserção de identidades) e
    extrapola para o limite de ruído zero.
  \item \textbf{PEC:} Aplica cancelamento probabilístico de erros
    através da decomposição de canais ruidosos em combinações
    convexas de operações ideais.
  \item \textbf{DD:} Insere pulsos de controle durante períodos de
    espera para desacoplar o sistema quântico do ambiente.
\end{enumerate}

A Tabela~\ref{tab:quantum} resume os principais artefatos do
subsistema quântico:

\begin{longtable}{p{4cm}p{3cm}p{5cm}}
  \caption{Artefatos do subsistema de computação quântica}
  \label{tab:quantum}\\
  \toprule
  \textbf{Componente} & \textbf{Arquivos} & \textbf{Descrição} \\
  \midrule
  \endfirsthead
  \toprule
  \textbf{Componente} & \textbf{Quantidade} & \textbf{Descrição} \\
  \midrule
  \endhead
  \bottomrule
  \endfoot
  Scripts Python/Rust & 26 & Implementações QML, MPS, Grad-CAM \\
  Citações acadêmicas & 21 & Referências Qualis A1 \\
  Saídas de validação & 7  & Resultados de execução e métricas \\
  Templates           & 24 & Estruturas de circuitos e pipelines \\
  Documentação        & 29 & Especificações, relatórios, diagramas \\
\end{longtable}

% ======================================================================
\section{ORQUESTRAÇÃO MULTIAGENTE NEXUS}
\label{sec:nexus}
% ======================================================================

O diretório \texttt{nexus/} contém 488 arquivos dedicados à
orquestração multiagente do ecossistema. A arquitetura Nexus
organiza-se em 6 camadas de meta-granularidade (L0--L6), cada uma
responsável por um nível de abstração na coordenação de agentes:

\begin{enumerate}
  \item \textbf{L0 -- Núcleo:} Tipos fundamentais, construtores de
    agentes e barreiras de sincronização primitivas.
  \item \textbf{L1 -- Comunicação:} Protocolos de troca de mensagens
    entre agentes, incluindo \textit{file-ipc} e \textit{fs-ipc}.
  \item \textbf{L2 -- Orquestração:} Coordenação de pipelines
    multiagente com barreiras de sincronização.
  \item \textbf{L3 -- Raciocínio:} Integração com os 4 motores de
    raciocínio formal (Z3, SymPy, Kanren, Critical).
  \item \textbf{L4 -- Aprendizado:} Loop evolutivo com registro de
    padrões de sucesso e geração de novas skills.
  \item \textbf{L5 -- Auditoria:} Verificação cruzada entre agentes,
    validação de consistência e geração de relatórios.
\end{enumerate}

O sistema conta com 120+ barreiras de sincronização e mais de 500
restrições de validação distribuídas entre as camadas. A orquestração
segue o padrão \textit{Barrier Synchronization}, onde agentes são
liberados para executar apenas quando todos os agentes do grupo
anterior completaram suas tarefas.

São 18 documentos de arquitetura de referência e 20 scripts Python
para automação de testes e validação do orquestrador.

% ======================================================================
\section{SUBSISTEMA DE CRIAÇÃO DE ARTIGOS ACADÊMICOS}
\label{sec:artigo}
% ======================================================================

O pipeline de criação de artigos acadêmicos, localizado em
\texttt{criador-artigo/}, compreende 91 arquivos organizados em
três subsistemas principais.

% ----------------------------------------------------------------------
\subsection{Agentes MASWOS (49 especialistas)}
% ----------------------------------------------------------------------

O sistema MASWOS (Multi-Agent Scientific Writing Orchestration System)
coordena 49 agentes especializados (identificados como agente-00 a
agente-44 mais 5 de suporte), cada um responsável por uma etapa do
processo de escrita acadêmica:

\begin{itemize}
  \item \textbf{Agentes de Pesquisa (00--09):} Busca, seleção e
    síntese de literatura científica.
  \item \textbf{Agentes de Estruturação (10--19):} Organização de
    seções, definição de hipóteses e metodologia.
  \item \textbf{Agentes de Redação (20--29):} Escrita de seções
    específicas (introdução, métodos, resultados, discussão).
  \item \textbf{Agentes de Revisão (30--39):} Revisão gramatical,
    de conteúdo e adequação a normas.
  \item \textbf{Agentes de Formatação (40--44):} Geração de LaTeX,
    PDF e metadados ABNT.
  \item \textbf{Agente Scheduler:} Coordenação temporal das etapas.
\end{itemize}

% ----------------------------------------------------------------------
\subsection{Pipeline de Qualidade}
% ----------------------------------------------------------------------

O pipeline de qualidade implementa múltiplos estágios de verificação:

\begin{itemize}
  \item \textbf{Revisão anti-AI:} Filtro TSAC com 87 palavras
    banidas para evitar detecção de escrita por IA.
  \item \textbf{Validação cruzada:} Correlação de Pearson entre
    seções, 3 níveis de consistência.
  \item \textbf{Correção iterativa:} 5 revisores simulados,
    4 consultores doutores, 6 motores de correção.
  \item \textbf{Avaliação automática:} Script
    AUTO\_SCORE\_QUALIS.py com 10 critérios e pesos de banca.
  \item \textbf{Corretor linguístico:} Detecção e remoção de
    caracteres CJK, verificação gramatical PT-BR.
\end{itemize}

O score mínimo para aprovação é 95/100, alinhado ao padrão Qualis A1.

% ----------------------------------------------------------------------
\subsection{Templates e Integração SDD--TDD}
% ----------------------------------------------------------------------

O diretório \texttt{criador-artigo/} contém 24 templates LaTeX, 14
referências Qualis A1 e integração completa com o pipeline
SDD--TDD--Harness--Hooks para validação de qualidade acadêmica.

% ======================================================================
\section{PIPELINE SDD--TDD--HARNESS--HOOKS E VALIDADORES}
\label{sec:validadores}
% ======================================================================

O pipeline de validação acadêmica segue o fluxo quadrifásico:
\[
\text{SDD} \longrightarrow \text{TDD} \longrightarrow
\text{Harness} \longrightarrow \text{Hooks}
\looparrowleft \text{Feedback}
\]

% ----------------------------------------------------------------------
\subsection{SDD -- Spec-Driven Development}
% ----------------------------------------------------------------------

O SDD estabelece que toda implementação deve ser precedida por uma
especificação formal. Cada especificação define metadados
(identificador, categoria, versão, autor), requisitos funcionais,
Assertivas Críticas (CAs) e dependências. O ecossistema mantém 186
componentes documentados com 100\% de cobertura.

% ----------------------------------------------------------------------
\subsection{TDD -- Test-Driven Development Acadêmico}
% ----------------------------------------------------------------------

O sistema TDDAcademic é composto por 5 especificações e 25 CAs:

\begin{longtable}{p{4cm}p{3cm}p{5cm}}
  \caption{Especificações TDDAcademic}
  \label{tab:tdd}\\
  \toprule
  \textbf{Especificação} & \textbf{ID} & \textbf{CAs} \\
  \midrule
  \endfirsthead
  \toprule
  \textbf{Especificação} & \textbf{ID} & \textbf{CAs} \\
  \midrule
  \endhead
  \bottomrule
  \endfoot
  TDDAcademic Base       & SPEC-TDDA-BASE      & 5 \\
  Ciclo de Desenvolvimento & SPEC-TDDA-CICLO   & 5 \\
  Baby Steps             & SPEC-TDDA-BABY      & 5 \\
  Regras de Refatoração  & SPEC-TDDA-REGRAS    & 5 \\
  Ferramentas TDD        & SPEC-TDDA-FERRAMENTAS & 5 \\
\end{longtable}

% ----------------------------------------------------------------------
\subsection{Arquitetura de Validadores}
% ----------------------------------------------------------------------

A arquitetura de validadores segue os princípios SOLID
\cite{martin2000principles} e padrões GoF \cite{gamma1995design}.
O núcleo é definido pela interface genérica \texttt{IValidator<T>}
com os métodos \texttt{validate(T)}, \texttt{canHandle(Spec)} e
\texttt{getCategory()}. A classe abstrata
\texttt{AbstractValidator<T>} implementa o padrão \textit{Template
Method} com o fluxo \texttt{preValidate $\rightarrow$ doValidate
$\rightarrow$ postValidate}.

As implementações concretas incluem \texttt{SpecValidator}
(SPEC-VALIDATOR-001), \texttt{CAValidator} (CA-VALIDATOR-001) e
\texttt{HarnessValidator} (HARNESS-VALIDATOR-001).

% ----------------------------------------------------------------------
\subsection{Hooks -- Automação Orientada a Eventos}
% ----------------------------------------------------------------------

O sistema de hooks implementa o padrão \textit{Observer}, permitindo
que agentes sejam notificados sobre eventos do pipeline:

\begin{itemize}
  \item \textbf{Pre-Hook (\texttt{on-spec-create}):} Valida estrutura
    e consistência antes da criação de especificações.
  \item \textbf{Error-Hook (\texttt{on-ca-fail}):} Registra falhas
    e notifica agentes responsáveis.
  \item \textbf{Post-Hook (\texttt{on-validation-complete}):} Gera
    relatórios e dispara implantações.
\end{itemize}

% ======================================================================
\section{SISTEMA EVOLUTIVO}
\label{sec:evolucao}
% ======================================================================

O sistema evolutivo do OpenCode (\texttt{evolution/}) implementa um
mecanismo de evolução autônoma inspirado em algoritmos genéticos. São
14 gerações documentadas, cada uma com score, insights extraídos e
novas skills geradas:

\begin{longtable}{p{2cm}p{4cm}p{2cm}p{4cm}}
  \caption{Evolução do ecossistema ao longo de 14 gerações}
  \label{tab:evolucao}\\
  \toprule
  \textbf{Round} & \textbf{Skill gerada} & \textbf{Score} & \textbf{Insight} \\
  \midrule
  \endfirsthead
  \toprule
  \textbf{Round} & \textbf{Skill gerada} & \textbf{Score} & \textbf{Insight} \\
  \midrule
  \endhead
  \bottomrule
  \endfoot
  1  & Validação quantitativa cruzada & 85  & Educação r=-0,03; P\&D r=+0,73 \\
  2  & Pipeline de artigos acadêmicos & 90  & Serviços de alta tecnologia r=+0,95 \\
  3  & Citações TSAC, Sci-Hub         & 92  & 46 anotações TSAC auditáveis \\
  4  & Loop de correção iterativa v2  & 95  & Revisores + consultores + corretores \\
  5  & Corretor CJK                   & 98  & Tolerância zero a vazamento CJK \\
  6  & Editais BR v2.0                & 92  & Busca real com curl + Firefox UA \\
  7  & Editais BR v7.1 cache          & 94  & 52 editais curados, 27 UFs \\
  8  & SDD+TDD acadêmico + Arguição   & 94  & Nota DAP 8,07$\rightarrow$9,0 \\
  9  & LaTeX refino + Framework Docs  & 96  & 4 overfulls eliminados, 16/16 TDD \\
  10 & Menu adaptativo + Plugin       & 96  & Menu estático $\rightarrow$ adaptativo \\
  11 & CORA-Eval Benchmark            & 97  & 150 tarefas, baseline 0,67 \\
  12 & Science Skills + MCP Expansion & 98  & 9 skills + 28 datasets + 4 MCPs \\
  13 & Reasoning Engines              & 96  & Z3 + SymPy + Kanren + Critical \\
  14 & Edição final e consolidação    & 98  & Score máximo sustentado \\
\end{longtable}

% ======================================================================
\subsection{Mecanismo de Evolução}
% ======================================================================

O mecanismo \texttt{autoevolve} \cite{OpenCode2025} segue o ciclo:

\begin{enumerate}
  \item \textbf{PLAN:} Identifica lacunas no ecossistema com base em
    métricas de uso e feedback.
  \item \textbf{ACT:} Gera nova skill ou modifica skill existente.
  \item \textbf{REFLECT:} Avalia impacto da mudança através de
    validação cruzada com CAs existentes.
  \item \textbf{EXTRACT:} Extrai padrões de sucesso para reuso.
  \item \textbf{EVOLVE:} Incorpora aprendizado no ecossistema.
\end{enumerate}

A \textit{skill} de evolução autônoma \texttt{Manus Evolve v1.0} é
implementada como plugin em \texttt{plugins/manus-evolve.ts} e
aprende automaticamente a partir de padrões de sucesso observados
durante execuções anteriores.

% ======================================================================
\section{MOTORES DE RACIOCÍNIO}
\label{sec:raciocinio}
% ======================================================================

O OpenCode integra 4 motores de raciocínio formal, cada um com um
propósito específico:

% ----------------------------------------------------------------------
\subsection{Z3 (Prova Formal)}
% ----------------------------------------------------------------------

O Z3 Theorem Prover (v4.16, Microsoft Research) é integrado para
verificação formal de propriedades de especificações. Utilizado para:

\begin{itemize}
  \item Validação de invariantes em especificações SDD.
  \item Verificação de consistência em Assertivas Críticas (CAs).
  \item Prova de propriedades de segurança em pipelines multiagente.
  \item Model checking de máquinas de estado do sistema Reversa.
\end{itemize}

% ----------------------------------------------------------------------
\subsection{SymPy (Matemática Simbólica)}
% ----------------------------------------------------------------------

O SymPy (v1.14) é utilizado para manipulação simbólica em:

\begin{itemize}
  \item Diferenciação e integração simbólica para validação de
    modelos matemáticos.
  \item Resolução de sistemas de equações em especificações
    científicas.
  \item Simplificação algébrica de expressões derivadas de CAs.
\end{itemize}

% ----------------------------------------------------------------------
\subsection{MiniKanren (Lógica Relacional)}
% ----------------------------------------------------------------------

Implementação do miniKanren para programação lógica relacional:

\begin{itemize}
  \item Definição de relações entre entidades do grafo de
    conhecimento do ecossistema.
  \item Consultas lógicas sobre dependências entre especificações.
  \item Inferência de relações não explicitamente declaradas.
\end{itemize}

% ----------------------------------------------------------------------
\subsection{Critical (Análise de Falácias)}
% ----------------------------------------------------------------------

O motor Critical implementa detecção de 15 tipos de falácias lógicas
e vieses cognitivos em argumentos gerados por agentes:

\begin{itemize}
  \item \textbf{Falácias formais:} Negação do antecedente, afirmação
    do consequente, silogismo disjuntivo inválido.
  \item \textbf{Falácias informais:} Apelo à autoridade, falsa
    dicotomia, generalização apressada, correlação-causalidade.
  \item \textbf{Vieses cognitivos:} Viés de confirmação, ancoragem,
    disponibilidade, excesso de confiança.
\end{itemize}

O sistema classifica 212+ tipos de raciocínio em 27 categorias,
incluindo 10 estratégias de teoria dos jogos aplicadas a cenários
de debate multiagente.

% ======================================================================
\section{SKILLS CIENTÍFICAS}
\label{sec:ciencia}
% ======================================================================

O OpenCode integra 38 \textit{skills} científicas que fornecem acesso
programático a bases de dados e ferramentas de pesquisa:

\begin{longtable}{p{3.5cm}p{3cm}p{5.5cm}}
  \caption{Skills científicas do ecossistema OpenCode}
  \label{tab:ciencia}\\
  \toprule
  \textbf{Skill} & \textbf{Tipo} & \textbf{Funcionalidade} \\
  \midrule
  \endfirsthead
  \toprule
  \textbf{Skill} & \textbf{Tipo} & \textbf{Funcionalidade} \\
  \midrule
  \endhead
  \bottomrule
  \endfoot
  AlphaFold       & Estrutura      & Predição de estruturas de proteínas \\
  PubMed          & Literatura     & Busca em literatura biomédica \\
  ChEMBL          & Farmacologia   & Moléculas bioativas e alvos \\
  UniProt         & Proteínas      & Metadados de proteínas \\
  ClinVar         & Genética       & Significância clínica de variantes \\
  FoldSeek        & Estrutura      & Busca estrutural 3D de proteínas \\
  PyMOL           & Visualização   & Renderização de estruturas moleculares \\
  OpenAlex        & Literatura     & Busca acadêmica multidisciplinar \\
  JASPAR          & Genômica       & Perfis de ligação de TFs \\
  ENCODE cCREs    & Genômica       & Elementos regulatórios cis \\
  QuickGO         & Ontologia      & Ontologia gênica GO \\
  HPA             & Proteínas      & Expressão proteica em tecidos \\
  NCBI Sequence   & Sequências     & Sequências de nucleotídeos/proteínas \\
  UCSC            & Genômica       & Conservação evolutiva e TFBS \\
  EMBL-EBI OLS    & Ontologia      & Serviço de ontologias biomédicas \\
  MSA             & Sequências     & Alinhamento múltiplo de sequências \\
  MMseqs2/BLAST   & Sequências     & Busca de similaridade \\
  AlphaGenome     & Genômica       & Efeitos de variantes não-codificantes \\
\end{longtable}

As skills científicas acessam 28 datasets adicionais, incluindo
gnomAD (variantes populacionais), GTEx (expressão em tecidos),
Ensembl (genômica comparativa), PDB (estruturas de proteínas) e
STRING (interações proteína-proteína).

% ======================================================================
\section{BENCHMARK CORA-EVAL}
\label{sec:coraeval}
% ======================================================================

O CORA-Eval (Cognitive ORchestrated Argumentation Evaluation) é um
benchmark abrangente para avaliação de raciocínio científico em
agentes de IA. Compreende 150 tarefas distribuídas em 10 dimensões
e 4 níveis de dificuldade (Básico, Intermediário, Avançado,
Pesquisa):

\begin{longtable}{p{3cm}p{2cm}p{2cm}p{5cm}}
  \caption{Dimensões do CORA-Eval}
  \label{tab:coraeval}\\
  \toprule
  \textbf{Dimensão} & \textbf{Nível} & \textbf{CTs} & \textbf{Descrição} \\
  \midrule
  \endfirsthead
  \toprule
  \textbf{Dimensão} & \textbf{Nível} & \textbf{CTs} & \textbf{Descrição} \\
  \midrule
  \endhead
  \bottomrule
  \endfoot
  D1 -- Matemática     & N1--N3 & 8  & Pitágoras, Gauss, Fibonacci, \
                                      Primos, MDC, Trigonometria \\
  D2 -- Física         & N1--N2 & 8  & MRU, Queda Livre, Hooke, \
                                      Pêndulo, Gases \\
  D3 -- Estatística    & N2--N3 & 9  & Teste t, ANOVA, Bootstrap, \
                                      MCMC, PCA \\
  D4 -- Química        & N1     & 3  & Balanceamento, Massa Molar, \
                                      Concentração \\
  D5 -- Biologia       & N1     & 3  & Transcrição, Tradução, GC \\
  D6 -- Geociências    & N1     & 3  & Rochas, Temperatura, Atmosfera \\
  D7 -- Código         & N3     & 6  & Syntax, Logic, Types, Security \\
  D8 -- Literatura     & N1     & 3  & Claims, Citations, Classification \\
  D9 -- Metodologia    & N2     & 8  & ANOVA, Cohen, Shapiro, \
                                      Fatorial 2³ \\
  D10 -- Interdisciplinar & N4 & 3  & Geometria Diferencial, \
                                      Teoria de Gauge \\
  D11 -- Long Horizon  & N1--N4 & 150 & Raciocínio em DAG \\
\end{longtable}

O benchmark utiliza um rastreador Python com persistência JSON,
Q-Score UCB1 para seleção adaptativa de tarefas e CORA-Score como
métrica composta ponderada por 7 verificadores (V1--V7). O baseline
atual é CORA-Score = 0,67.

% ======================================================================
\section{INTEGRAÇÃO MIROFISH E BETTAFISH}
\label{sec:mirofish}
% ======================================================================

O ecossistema OpenCode integra 11 componentes dos projetos
MiroFish e BettaFish (666ghj/MiroFish, 666ghj/BettaFish):

\begin{itemize}
  \item \textbf{OASIS:} Perfil de agentes a partir de grafos de
    conhecimento (OASIS Profile Generator).
  \item \textbf{Forum:} Motor de debate multiagente com moderação
    LLM (ForumEngine).
  \item \textbf{Config:} Geração de configurações complexas com
    fallback heurístico (SimulationConfigGenerator).
  \item \textbf{Graph:} Construção e consulta de grafos de
    conhecimento (GraphBuilderService, GraphTools).
  \item \textbf{Report:} Geração de relatórios com cadeia ReACT
    (ReportAgent).
  \item \textbf{Nash:} Equilíbrio de Nash para payoff matrix em
    debates multiagente.
  \item \textbf{Stats:} Rigor estatístico com correções de
    Bonferroni e Cohen's d.
  \item \textbf{Qualis:} Auditoria Qualis A1 com verificação de
    conformidade.
  \item \textbf{Sensitivity:} Análise de sensibilidade de
    hiperparâmetros.
  \item \textbf{IMRAD:} Formatação de artigos no padrão Qualis A1.
  \item \textbf{Debate:} Estratégias de debate com 212+ tipos de
    raciocínio.
\end{itemize}

% ======================================================================
\section{SISTEMA DE PESQUISA SEER}
\label{sec:seer}
% ======================================================================

O SEER (Seeker for Evidence-based Enhanced Research) é um sistema
de pesquisa acadêmica composto por 78 arquivos e 10 agentes Python:

\begin{itemize}
  \item \textbf{Agente Searcher:} Consulta a 10+ fontes acadêmicas
    (arXiv, OpenAlex, Semantic Scholar, PubMed, CORE).
  \item \textbf{Agente Grounder:} Valida afirmações com evidências
    citáveis.
  \item \textbf{Agente Argument Tree:} Constrói árvores de argumentos
    com suporte empírico.
  \item \textbf{Agente Verifier:} Verifica cruzadamente afirmações
    entre múltiplas fontes.
\end{itemize}

O SEER suporta pesquisa aprofundada com rastreamento de confiança
por afirmação, produzindo relatórios em Markdown com citações
verificáveis.

% ======================================================================
\section{IMPLEMENTAÇÃO LGPD}
\label{sec:lgpd}
% ======================================================================

A implementação da Lei Geral de Proteção de Dados (LGPD) no
OpenCode é organizada em 3 fases:

\begin{itemize}
  \item \textbf{Fase 1 -- Análise (\texttt{fase1\_analise/}):}
    Mapeamento de dados pessoais, identificação de bases legais,
    análise de riscos e documentação de conformidade.
  \item \textbf{Fase 2 -- Implementação (\texttt{fase2\_implementacao/}):}
    Implementação de controles técnicos, anonimização, pseudonimização,
    registro de consentimento e DPO as a Service.
  \item \textbf{Fase 3 -- Governança:} Monitoramento contínuo,
    auditoria de conformidade, resposta a incidentes e relatórios
    de impacto.
\end{itemize}

% ======================================================================
\section{PLUGINS E COMANDOS}
\label{sec:plugins}
% ======================================================================

O ecossistema dispõe de 15 plugins (10 npm + 2 TypeScript locais +
3 bridge), 14 comandos no formato \texttt{/comando} e 1 LSP
(\textit{Language Server Protocol}) TypeScript:

\begin{longtable}{p{3cm}p{4cm}p{5cm}}
  \caption{Comandos do ecossistema OpenCode}
  \label{tab:comandos}\\
  \toprule
  \textbf{Comando} & \textbf{Função} & \textbf{MCPs/Plugins associados} \\
  \midrule
  \endfirsthead
  \toprule
  \textbf{Comando} & \textbf{Função} & \textbf{MCPs/Plugins associados} \\
  \midrule
  \endhead
  \bottomrule
  \endfoot
  /evolve  & Auto-evolução do ecossistema & autoevolve, ecosystem-sync \\
  /reversa & Engenharia reversa           & reversa-* agents, filesystem \\
  /plan    & Planejamento de escrita      & writing-plans, sequential-thinking \\
  /auto    & Automação geral              & openagent, todos MCPs \\
  /quantum & Pipeline quântico            & quantum-nexus-phd, code-runner \\
  /artigo  & Criação de artigos            & SEER + Article Creator \\
  /forum   & Debate multiagente            & agent-forum, Cora-Debate \\
  +7 comandos adicionais \\
\end{longtable}

% ======================================================================
\section{CONSIDERAÇÕES FINAIS}
\label{sec:conclusao}
% ======================================================================

O ecossistema OpenCode v5.0.0 representa uma contribuição original
para a convergência entre engenharia de software, inteligência
artificial, computação quântica e produção acadêmica. Com mais de
600 componentes integrados, 125 agentes especializados e 150
habilidades, a plataforma demonstra a viabilidade de um ambiente
unificado de desenvolvimento assistido por IA com rigor acadêmico.

As principais contribuições incluem:

\begin{enumerate}
  \item \textbf{Arquitetura em 3 camadas} com desacoplamento entre
    MCPs, Skills e Agentes, permitindo substituição independente
    de implementações.
  \item \textbf{Computação quântica aplicada} com QML atingindo
    89,52\% de acurácia em classificação de imagens médicas e
    pipeline completo de mitigação de erros.
  \item \textbf{Orquestração multiagente} com 488 arquivos Nexus,
    6 camadas de meta-granularidade e 120+ barreiras de
    sincronização.
  \item \textbf{Pipeline acadêmico MASWOS} com 49 agentes e score
    Qualis A1 sustentado acima de 95/100.
  \item \textbf{Sistema evolutivo autônomo} que elevou o score do
    ecossistema de 85 para 98 ao longo de 14 gerações.
  \item \textbf{Motores de raciocínio formal} integrados (Z4,
    SymPy, miniKanren, Critical) com 212+ tipos em 27 categorias.
  \item \textbf{Benchmark CORA-Eval} com 150 tarefas em 10
    dimensões acadêmicas.
  \item \textbf{38 skills científicas} integradas a bases como
    AlphaFold, PubMed, ChEMBL e UniProt.
\end{enumerate}

Entre as direções para trabalho futuro, destacam-se a expansão do
sistema de computação quântica para 100+ qubits, a integração com
verificadores formais adicionais (Coq, Lean 4), a ampliação do
CORA-Eval para 500+ tarefas e a publicação do ecossistema como
plataforma open-source para validação pela comunidade acadêmica.

% ======================================================================
\bibliography{raio_x_opencode}
% ======================================================================

\end{document}
